\subsection{The canonical projectivity ideal}
\label{sec:canonical-ideal}
We prove Theorem~\ref{thm:canonical-ideal}. This removes the choice of a
support section in the preceding criterion. Put $\mathcal A=\mathcal A_L$.
For a $Q$-linear $\phi:L\to Q$, the groups
\[
 N_\phi(k)=F(\phi(k))
\]
form an $\mathcal A$-module by pulling back the rank-one free module.
There is no assertion that $N_\phi$ is projective. For an admissible
pair $(l,\phi)$ there is a natural map to a representable,
\[
 \theta_{l,\phi}:N_\phi\longrightarrow\mathcal A(l,-),\qquad
 x\longmapsto t_{\phi(k)l,k}a_{\phi(k),x,l}
 \quad(x\in N_\phi(k)).
\]
The transition and scalar naturality relations defining $\mathcal A$
prove its naturality.

\begin{proof}[Proof of Theorem~\ref{thm:canonical-ideal}]
Postcomposition preserves the displayed spans by naturality of
$\theta_{l,\phi}$. Precomposition by $t_{h,l}$, for $h\leq l$, uses the
same $\phi$, now admissible at $h$. The resulting arrow is
$\theta_{h,\phi}(k)(x)$.
For precomposition by $a_{q,u,h}$ with $qh=l$, put
$\psi(k)=\phi(k)q$. Right multiplication by the fixed $q$ makes $\psi$
a left $Q$-linear map even when $Q$ is noncommutative. Moreover
$\psi(k)h=\phi(k)l\leq k$. Scalar composition gives
\[
 \theta_{l,\phi}(k)(x)a_{q,u,h}
 =\theta_{h,\psi}(k)\bigl(\mu_{\phi(k),q}(x,u)\bigr).
\]
Thus precomposition also preserves the spans. Additivity and composition
of generators prove that $\mathcal J_L$ is already a two-sided ideal.
This includes $F(0)\ne0$: the zero functional is admissible and its
coefficient group is retained.

Suppose $M$ is a retract of $T^n$, with support maps
$\sigma:L\to Q^n$, $\rho:Q^n\to L$, and $\rho\sigma=1$.
Let $l_i=\rho(e_i)$ and $\phi_i=\sigma_i$. Then
\[
 k=\bigvee_i\phi_i(k)l_i,
\]
so every $(l_i,\phi_i)$ is admissible. Let $m_i\in G(l_i)$ be the image
of the $i$th free generator. On amplitudes the free retraction is the
sum of the evaluation maps
\[
 N_{\phi_i}\xrightarrow{\theta_{l_i,\phi_i}}
 \mathcal A(l_i,-)\xrightarrow{m_i}G.
\]
Its composite with the section is $1_G$, and every such evaluation has
image in $\mathcal J_LG$. Hence $\mathcal J_LG=G$. The support is a
finite free retract. Finite generation of $G$ follows from
Proposition~\ref{prop:ringoid-finite}; its ordinary projectivity follows
by lifting a pointwise amplitude epimorphism with identity support, as in
the proof of Theorem~\ref{thm:ringoid-projective}.

Conversely suppose all three conditions hold. Choose finitely many
$\mathcal A$-generators of $G$. The equality $\mathcal J_LG=G$ expresses
each as a finite sum of values of the defining arrows on elements
$m_i\in G(l_i)$. Collect the finitely many admissible pairs and elements
occurring in these expressions. Their evaluation maps give a pointwise
epimorphism
\[
 c:\bigoplus_{i=1}^nN_{\phi_i}\longrightarrow G.
\]
Indeed its image is a submodule containing all the chosen generators.
Projectivity of $G$ supplies an $\mathcal A$-linear section $s$.
The corresponding support maps satisfy only
\[
 \rho_1\phi(k)=\bigvee_i\phi_i(k)l_i\leq k.
\]
Choose independently a finite free support retraction
$\sigma_0:L\to Q^m$, $\rho_0:Q^m\to L$. Appending its coordinates gives
\[
 \sigma'=(\phi,\sigma_0),\qquad
 \rho'(u,v)=\rho_1(u)\vee\rho_0(v),\qquad
 \rho'\sigma'=1_L.
\]
Map the first $n$ free generators of $T^{n+m}$ to the $m_i$, and the
remaining generators to the group zeros over $\rho_0(e_j)$.
These zeros have the required supports although their amplitudes vanish.
The resulting $T$-map has support $\rho'$. Give its section support
$\sigma'$ and amplitude $(s,0)$. The fixed-support equivalence shows
that it is $T$-linear. Its composite has identity support and amplitude
$cs=1_G$. Hence $M$ is a finite free retract.

Finally, extension along the quotient category is the pointwise quotient
$G/\mathcal J_LG$: its universal property is precisely that maps annihilate
the action of $\mathcal J_L$. This proves the tensor formulation.
The fixed-support equivalence gives the claimed classification there.
A $Q$-isomorphism of supports transports admissible functionals by
precomposition and transports scalar and transition arrows, hence the
ideal and all the stated data.
\end{proof}

\begin{corollary}[idempotents in the canonical ideal]
\label{cor:canonical-ideal-matrices}
The amplitudes in Theorem~\ref{thm:canonical-ideal} are exactly the images
of idempotent endomorphisms $E$ of finite sums
$P=\bigoplus_i\mathcal A_L(l_i,-)$ whose matrix entries lie in
$\mathcal J_L$. Invariantly, the condition is $E(P)\subseteq\mathcal J_LP$.
The usual rectangular intertwiners classify these images up to isomorphism.
\end{corollary}
\begin{proof}
At an object $k$, the submodule $\mathcal J_LP$ is
$\bigoplus_i\mathcal J_L(l_i,k)$, by the ideal property and the identity
generators of the representables. If $G=\im E$ and $\mathcal J_LG=G$,
its inclusion into $P$ gives $E(P)\subseteq\mathcal J_LP$, exactly the
entry condition. Conversely, if $g\in G\subseteq\mathcal J_LP$, write
$g=\sum_a j_ap_a$ and apply $E$. Then
$g=\sum_a j_aE(p_a)\in\mathcal J_LG$.
Every finite projective has such a representable presentation.
Extending an isomorphism of images by inclusion and projection gives
the rectangular intertwiners, and restricting those identities recovers
the isomorphism.
\end{proof}
The matrix argument is the ordinary trace-ideal argument; compare
Herbera and P\v r\'ihoda \cite[p.~2 of the preprint]{HerberaPrihoda}.
The additional datum here is the explicitly defined ideal
$\mathcal J_L$. For the top-born diagram in
Proposition~\ref{prop:ringoid-obstruction}, admissibility at the top
forces $\phi(a)=\phi(b)=0$, hence $\phi(a\vee b)=0$; therefore
$\mathcal J_LG=0\ne G$. The canonical ideal detects the same obstruction.

There is no claim that these categories are left rooted. The projective
functor classification of Cort\'es-Izurdiaga and Torrecillas
\cite[Theorem~1, pp.~320--322]{CortesTorrecillas} assumes that hypothesis;
their Remark~3 and Example~1 describe the cycle issue. The action
categories here can have scalar loops, so that result does not replace
ordinary projectivity in Theorem~\ref{thm:canonical-ideal}. Nor does the
theorem give a finite procedure for listing all admissible functionals.
